% DANH MỤC KÝ HIỆU VÀ CHỮ VIẾT TẮT: đặt sau Mục lục, trước Chương 1; mọi ký hiệu hoặc chữ viết tắt mới xuất hiện trong báo cáo được thêm vào bảng này.
\clearpage
\phantomsection
\section*{DANH MỤC KÝ HIỆU VÀ CHỮ VIẾT TẮT}
\addcontentsline{toc}{section}{Danh mục ký hiệu và chữ viết tắt}

\begin{center}
\small
\renewcommand{\arraystretch}{1.2}
\begin{tabular}{|p{2.6cm}|p{14.4cm}|}
\toprule
Ký hiệu & Ý nghĩa \\
\midrule
\multicolumn{2}{|l|}{\textit{Bộ dữ liệu}} \\
HU & Hồ Huron: mực nước hồ hằng năm \\
LY & Linh miêu: logarit số linh miêu bị bẫy hằng năm \\
DR & DAX dạng suất sinh lợi: suất sinh lợi log hằng ngày của chỉ số DAX \\
NI & Nile: lưu lượng hằng năm của sông Nile \\
AP & AirPassengers: logarit số hành khách hàng không hằng tháng \\
DP & DAX dạng log giá: logarit giá đóng cửa hằng ngày của chỉ số DAX \\
DAX, SMI & chỉ số chứng khoán của Đức, chỉ số chứng khoán của Thụy Sĩ \\
DL1--DL6 & nguồn dữ liệu đánh số trong mục Tài liệu tham khảo \\
\midrule
\multicolumn{2}{|l|}{\textit{Mô hình}} \\
AR & tự hồi quy (autoregressive) \\
MA & trung bình trượt (moving average) \\
ARMA & tự hồi quy trung bình trượt \\
ARIMA & ARMA cho chuỗi đã lấy sai phân (tích hợp) \\
SARIMA & ARIMA có thành phần mùa vụ \\
GARCH & mô hình phương sai có điều kiện thay đổi theo thời gian \\
$I(d)$ & chuỗi tích hợp bậc $d$, tức chuỗi dừng sau $d$ lần lấy sai phân \\
WN, iid & nhiễu trắng; độc lập và cùng phân phối \\
\midrule
\multicolumn{2}{|l|}{\textit{Phương pháp ước lượng và kiểm định}} \\
ACF, PACF & hàm tự tương quan, hàm tự tương quan riêng phần \\
YW, OLS, ML & ước lượng Yule--Walker, bình phương tối thiểu, cực đại hợp lý \\
LB, $Q_{LB}$ & kiểm định Ljung--Box và thống kê của nó \\
ADF, KPSS & kiểm định nghiệm đơn vị Dickey--Fuller mở rộng; kiểm định tính dừng Kwiatkowski--Phillips--Schmidt--Shin \\
DM & kiểm định Diebold--Mariano \\
AIC, BIC & tiêu chuẩn thông tin Akaike, tiêu chuẩn thông tin Bayes \\
SD, SE & độ lệch chuẩn; sai số chuẩn \\
\midrule
\multicolumn{2}{|l|}{\textit{Thước đo sai số}} \\
MAE, RMSE & sai số tuyệt đối trung bình; căn bậc hai của sai số bình phương trung bình \\
MAPE, MASE & sai số phần trăm tuyệt đối trung bình; sai số tuyệt đối trung bình đã chuẩn hóa \\
MCSE & sai số Monte Carlo của một kết quả mô phỏng \\
\midrule
\multicolumn{2}{|l|}{\textit{Quy ước trong báo cáo}} \\
CH1--CH3 & ba câu hỏi nghiên cứu nêu ở \muc{1.2} \\
P, M, K & xếp loại phù hợp, phù hợp một phần, không phù hợp \\
\bottomrule
\end{tabular}
\end{center}
